\begin{frame}{발표 순서}
  \small
  \begin{columns}[T]
    \begin{column}{0.5\textwidth}
      \begin{itemize}\setlength\itemsep{4pt}
        \item[Q1] Strict AOT — 재사용 vs 재설계
        \item[Q2] 고정 shape 버킷 / warm-up의 한계
        \item[Q3] Prefix caching과 prefill 버킷 설계
        \item[Q4] 포팅 계획이 틀어진 사례
        \item[Q5] 컴파일 스펙트럼별 강제사항
      \end{itemize}
    \end{column}
    \begin{column}{0.5\textwidth}
      \begin{itemize}\setlength\itemsep{4pt}
        \item[Q6] RadixAttention vs PagedAttention
        \item[Q7] 스케줄러 차이 · host 병목
        \item[Q8] RadixAttention의 RBLN 포팅 경계
        \item[Q9] AOT에서 speculative decoding
        \item[Q10] TPOT P99 악화 분석
      \end{itemize}
    \end{column}
  \end{columns}
\end{frame}

\begin{frame}{전제 — vllm-rbln에서 읽어낸 RBLN 스택}
  \begin{itemize}\setlength\itemsep{5pt}
    \item \textbf{완전 AOT} — 런타임 컴파일 / eager fallback 없음; 미준비 shape은 실행 불가.
    \item \textbf{두 경로} — optimum-rbln \texttt{.rbln} 아티팩트 \,/\, \texttt{torch.compile(dynamic=False)} + rebel backend.
    \item \textbf{Prefill 단독}(batch $=1$); \textbf{decode batch만 버킷팅}; seq len은 고정 \texttt{max\_model\_len} 그래프.
    \item \textbf{KV block $=$ attention partition} (1--4K tok) $\Rightarrow$ sub-block prefix caching; hit 시 KV \emph{복사}(동기).
    \item \textbf{TP는 그래프에 컴파일}; \textbf{overlap scheduler 없음}.
  \end{itemize}
  \begin{alertblock}{한 줄}
    RBLN은 strict-AOT — 미준비 shape에 런타임 탈출구가 없다. 나머지 설계는 여기서 파생.
  \end{alertblock}
\end{frame}
